\section{Conclusion}\label{conclusion}

We present a broad model--representation LOGO benchmark showing that regularised linear models and auditable tree ensembles
using physics-motivated signal representations achieve competitive and often lower
point-estimate errors than the evaluated deep-learning baselines under
strict leave-one-condition-out validation. The main conclusions are:
\begin{itemize}
\tightlist
\item
  Random forests on pass-level fused acoustic-emission and vibration
	  dB-z and log-mel descriptors form the leading observed point-estimate
	  group. The reproducible current fixed dB-z pipeline gives
	  0.0210~\(\upmu\)m, and cache-matched inner-grouped reruns give
	  0.0207~\(\upmu\)m for dB-z and 0.0206~\(\upmu\)m for log-mel.
	  The nearly equal means mask different tails: nested maximum fold MAE is
	  0.1402~\(\upmu\)m for dB-z and 0.2061~\(\upmu\)m for log-mel, making
	  dB-z the better observed worst-condition candidate in this dataset.
	  Historical canonical values of 0.0200 and 0.0201~\(\upmu\)m are retained
	  for provenance and the planned post-selection comparison, not as a
	  uniquely reproducible winner. The primary planned comparison of 14
  Holm--Bonferroni-corrected Wilcoxon tests finds clear superiority only
  over the implementation-specific historical shallow MLP baseline; a corrected
  grouped-validation MLP sensitivity remains worse, while the remaining pairwise differences are
  favourable by point estimate but not significant. This family is an
  artifact-based post-selection analysis of preserved historical predictions,
  not an independently regenerated comparison. A secondary current-protocol
  five-seed analysis produces the same qualitative ordering for the selected
  repeated baselines, although it does not exactly reproduce the historical
  tuned checkpoints. The ranking among the
  top RF variants is exploratory because the MAE differences are smaller
  than the available measurement-resolution evidence, and the comparison
	  remains non-nested/post-selection.
\item
	  In the matched 14-condition modality ablation, vibration alone gives
	  0.0217~\(\upmu\)m versus 0.0210~\(\upmu\)m for fusion. Vibration therefore
	  carries most of the predictive signal in the present setup, while the
	  AE mean-level gain is below the practical metrology scale.
\item
  Quantitative explainability analysis associates model decisions with
  physically plausible frequencies: vibration importance lies in the
  chatter/structural and high-frequency bands, and acoustic-emission
  importance has leading peaks near 334 and 883~kHz, with 34.21\% of the
  current RF's out-of-fold AE attribution above the nominal 1~MHz response
  limit; this requires sensor-transfer-function verification before it
  can be treated as validated physics.
\item
  Uncertainty calibration is currently the weakest component. Full
  16-fold LOGO MC-dropout on the separately regenerated ResNetVibCNN
  checkpoint set (deterministic MAE 0.0443~\(\upmu\)m) gives only 50.2\% empirical coverage
  (0\% on Condition~7), and OOB-calibrated tree-standard-deviation
  intervals for a separate 15-condition RF uncertainty variant under-cover the nominal 95\% level
  (80.3\% sample-weighted coverage, 0\% on Condition~7). Future
  calibration must be condition-aware.
\item
  The fitted RF depends predictively on vibration content above 15~kHz, but
  this region lies close to the 25.6~kHz Nyquist limit. Without archived
  accelerometer, anti-alias-filter, and mounting transfer functions, this
  dependence cannot be assigned to a mechanical mechanism.
\item
  Deployment timing is limited to feature-precomputed model-inference
  latency. The cached descriptors are full-pass aggregates, so the present
  benchmark does not establish within-pass, closed-loop, or end-to-end
  prediction latency; streaming aggregation and independent external
  validation remain future work.
\end{itemize}

For engineering-informatics practice, the results support considering
physics-motivated signal representations and auditable classical learners as strong baseline
models that should be evaluated before selecting higher-capacity deep
architectures. On small-data monitoring
problems, this combination can deliver competitive point-estimate
accuracy and auditability through tree-based attribution methods; the
evidence is not strong enough to claim universal superiority over deep
learning and full online latency remains future work. These findings
remain limited to the present machine, wheel, material, and sensor setup
and require external validation before industrial generalisation. Future
work will archive all profilometer traces, compute gauge R\&R, add
cyclostationary and multi-window features, extend uncertainty calibration
to all top models, report end-to-end latency from raw signals, and
validate latency-memory trade-offs on dedicated edge hardware.
